\let\negmedspace\undefined
\let\negthickspace\undefined
\documentclass[journal]{IEEEtran}
\usepackage[a5paper, margin=10mm, onecolumn]{geometry}
%\usepackage{lmodern} % Ensure lmodern is loaded for pdflatex
\usepackage{tfrupee} % Include tfrupee package

\setlength{\headheight}{1cm} % Set the height of the header box
\setlength{\headsep}{0mm}     % Set the distance between the header box and the top of the text

\usepackage{gvv-book}
\usepackage{gvv}
\usepackage{cite}
\usepackage{physics}
\usepackage{amsmath,amssymb,amsfonts,amsthm}
\usepackage{algorithmic}
\usepackage{graphicx}
\usepackage{textcomp}
\usepackage{xcolor}
\usepackage{txfonts}
\usepackage{listings}
\usepackage{enumitem}
\usepackage{mathtools}
\usepackage{gensymb}
\usepackage{comment}
\usepackage[breaklinks=true]{hyperref}
\usepackage{tkz-euclide} 
\usepackage{listings}
% \usepackage{gvv}                                        
\def\inputGnumericTable{}                                 
\usepackage[latin1]{inputenc}                                
\usepackage{color}                                            
\usepackage{array}                                            
\usepackage{longtable}                                       
\usepackage{calc}                                             
\usepackage{multirow}                                         
\usepackage{hhline}                                           
\usepackage{ifthen}                                           
\usepackage{lscape}
\begin{document}

\bibliographystyle{IEEEtran}

\title{10.7.71}
\author{EE25BTECH11022 - sankeerthan}
% \maketitle
% \newpage
% \bigskip
{\let\newpage\relax\maketitle}

\renewcommand{\thefigure}{\theenumi}
\renewcommand{\thetable}{\theenumi}
\setlength{\intextsep}{10pt} % Space between text and floats


\numberwithin{align}{enumi}
\numberwithin{figure}{enumi}
\renewcommand{\thetable}{\theenumi}

\textbf{Question:}The area formed by the positive X axis and the normal and the tangent to the circle$x^2 + y^2 = 4$ at $\brak{1,\sqrt{3}}$ is \underline{\phantom{imagine}}. 


\textbf{Solution:}
Equation of the circle

\begin{align}
\vec{X}^T \vec{X} = 4 ,\text{where} \vec{X} =\myvec{x \\ y}
\end{align}

Given point of contact

\begin{align}
\vec{P} = \myvec{1 \\ \sqrt{3}}
\end{align}

Equation of the tangent at $\vec{P}$

\begin{align}
\vec{P}^T \vec{X} = 4
\end{align}

Let intersection of tangent with X-axis be $\vec{A} =\myvec{a \\ 0}$

\begin{align}
\vec{P}^T \vec{A} &= 4  \\
\myvec{1 & \sqrt{3}}\myvec{a \\ 0} &= 4 \implies a = 4 \\
\vec{A} = \myvec{4 \\ 0}
\end{align}

Equation of the normal at $\vec{P}$

\begin{align}
\vec{X} = t\myvec{1 \\ \sqrt{3}},t \in \mathbb{R}
\end{align}

Let intersection of normal with X-axis be $\vec{N} = \myvec{x \\ 0}$ 

\begin{align}
\vec{N}= t\myvec{1 \\ \sqrt{3}}
0 = t\sqrt{3} \implies t = 0 \implies \vec{N} = \myvec{0 \\ 0}
\end{align}

Area of triangle formed by points $\vec{O}$,$\vec{A}$,$\vec{P}$

\begin{align}
\vec{O} =\myvec{0 \\ 0},\vec{A} =\myvec{4 \\ 0},\vec{P} = \myvec{1 \\ \sqrt{3}} 
   \end{align} 

   \begin{align}
S &=\frac{1}{2}\left|\myvec{0 & 0 & 1 \\4 & 0 & 1 \\1 & \sqrt{3} & 1} \right| \\
&=\frac{1}{2}|0\brak{0-\sqrt{3}}-0\brak{4-1}+1\brak{4\sqrt{3}-0}| \\
&=\frac{1}{2}|4\sqrt{3}|=2\sqrt{3} 
\end{align}

Therefore, required area formed =$2\sqrt{3}$


\begin{figure}[h!]
    \centering
    \includegraphics[width=1\columnwidth]{figs/points.png}
    \caption{}
    \label{}
\end{figure}


\end{document}
